\section{Análise no Tempo Contínuo}

\subsection{Comportamento Dinâmico do Sistema}

A análise inicial do sistema foi realizada por meio da aplicação de perturbações em malha aberta nas variáveis de entrada.
As respostas temporais foram obtidas por integração numérica do modelo não linear utilizando a biblioteca \textit{SciPy} para Python \cite{virtanen_2020}.
O objetivo dessa etapa é identificar qualitativamente a influência de cada entrada sobre os estados do processo.

A Figura~\ref{fig:step_fr} apresenta a resposta temporal dos estados para uma perturbação do tipo degrau aplicada à vazão de reciclo $F_R$.
Observa-se que essa alteração provoca, em regime permanente, uma elevação nas temperaturas do sistema e um aumento na fração molar do componente $B$ nas três unidades do processo.
Esse comportamento é fisicamente coerente, pois a recirculação de reagente não convertido aumenta a formação de $B$ e intensifica a taxa global de reação, elevando as temperaturas em todas as unidades.

\begin{figure}[ht]
  \centering
  \includegraphics[width=0.8\textwidth]{figuras/step-FR-outputs.png}
  \caption{Resposta temporal dos estados do sistema a uma perturbação do tipo degrau aplicada à vazão de reciclo $F_R$, cujo valor é alterado de 17 $\mathrm{m^3 \cdot h^{-1}}$ para 25,5 $\mathrm{m^3 \cdot h^{-1}}$ em $t =$ 0,2 h.}
  \label{fig:step_fr}
\end{figure}

Para sintetizar os efeitos de todas as entradas sobre os estados do sistema, a Tabela~\ref{tab:interacoes}, apresentada no Apêndice~\ref{sec:apendice_interacoes}, reúne qualitativamente os efeitos observados em cada estado após a aplicação de perturbações individuais.

\subsection{Linearização}

O modelo não linear do sistema reator-separador foi linearizado por meio de expansão em série de Taylor de primeira ordem em torno do ponto de operação em regime estacionário apresentado no Apêndice~\ref{sec:apendice_ponto_operacao}.
A linearização tem como objetivo viabilizar a análise das interações entre entradas e saídas por meio de funções de transferência, bem como o estudo da estabilidade e da resposta dinâmica do sistema no domínio linear.

A validade da aproximação linear é restrita à vizinhança do ponto de operação, sendo adequada apenas para pequenas perturbações em torno desse regime.
A Figura~\ref{fig:linear_vs_nonlinear}, apresentada no Apêndice~\ref{sec:apendice_linearizacao}, ilustra esse comportamento ao comparar a resposta do modelo linearizado e do modelo não linear.
Observa-se boa concordância entre os modelos apenas para variações próximas ao ponto de operação.
Em particular, nota-se que a variável $x_{B1}$ assume valores superiores a 1 no modelo linearizado, o que não possui significado físico.

\subsection{Obtenção das Funções de Transferência}

O diagrama de blocos do sistema reator-separador, apresentado no Apêndice~\ref{sec:apendice_diagrama_blocos}, evidencia o forte acoplamento entre as variáveis do processo.
Nesse cenário, a obtenção das funções de transferência traduz essa rede interconectada em relações diretas de entrada e saída, facilitando sua análise.

As funções de transferência foram obtidas no domínio de Laplace, a partir do modelo linearizado, por meio da relação matricial:
\begin{equation}
  G(s) = (sI - A)^{-1}B,
\end{equation}
em que $A$ e $B$ são, respectivamente, as matrizes de estado e de entrada do modelo linearizado em torno do ponto de operação, $I$ é a matriz identidade e $s$ é a variável de Laplace.
A matriz $G(s)$ representa a matriz de funções de transferência do sistema, na qual cada elemento $g_{i,j}(s)$ descreve a relação dinâmica entre a entrada $j$ e a saída $i$.

Em seguida, foi aplicado um procedimento de simplificação visando tornar a representação mais compacta e adequada à análise, incluindo o cancelamento de polos e zeros coincidentes, a remoção de dinâmicas de alta frequência pouco relevantes e o ajuste do ganho estático para preservação do comportamento em regime permanente.
Todos os elementos $g_{i,j}(s)$ resultantes são apresentados no Apêndice~\ref{sec:apendice_gs}.

\subsection{Análise do canal $T_1(s)/F_R(s)$}

Como exemplo, considera-se o canal que relaciona a temperatura do primeiro reator $T_1$ à vazão de reciclo $F_R$.
A função de transferência correspondente é a $g_{1,3}(s)$, dada pela Equação~\ref{eq:g13} no Apêndice~\ref{sec:apendice_gs}.

O ganho estático positivo desse canal indica que um aumento na vazão de reciclo resulta, em regime permanente, em um aumento da temperatura do primeiro reator. Além disso, o diagrama de polos e zeros apresentado no Apêndice~\ref{sec:apendice_pzmap} mostra que todos os polos possuem parte real negativa, indicando que o canal é estável. Observa-se também a presença de um zero no semiplano direito, característica associada à ocorrência de resposta inversa.

A resposta temporal, em desvio, foi obtida a partir da aplicação da transformada inversa de Laplace à função de transferência, considerando uma entrada degrau unitário. Dessa forma, a solução analítica é expressa por:
\begin{equation}
  y(t)=0{,}267 + 0{,}12 e^{- 48{,}038 t} \sin{\left(19{,}62 t \right)} + 0{,}071 e^{- 48{,}038 t} \cos{\left(19{,}62 t \right)} - 0{,}338 e^{- 3{,}179 t}.
  \label{eq:step_g13}
\end{equation}

A comparação entre essa solução e a resposta simulada por meio da biblioteca \textit{Control} para Python \cite{python-control2021} é apresentada no Apêndice~\ref{sec:apendice_step_g13}, evidenciando a concordância entre ambas. Embora os termos senoidais sejam rapidamente amortecidos pelo fator exponencial $e^{-48{,}038t}$, a presença de um zero no semiplano direito torna inadequada a aproximação por um modelo de primeira ordem, pois sua eliminação compromete a representação da resposta transitória. Nesse contexto, uma aproximação que substitua apenas o par de polos complexos conjugados por um polo real equivalente pode ser mais viável para preservar as principais características da resposta dinâmica do sistema.
